% This file should be rename with the speaker's last name. (e.g : smith)

% Write the author names in the order you wish them to appear
% and underline the one who will give the talk.

% Only the speaker's email address is required.

\begin{abstract_online}{Investigations with DGS and CAS dealing with problems of equal area and particularly a possible generalization to 3D of the known results of the Lhuillier problem}{%
        \underline{Jean-Jacques Dahan}$^{1}$}{[jjdahan@wanadoo.fr]
        }{%
        $^1$ Head of the research group of dynamic geometry, IRES of Toulouse, Paul Sabatier University.}
    
This presentation aims to illustrate the dialectic between DGS and CAS during investigations which goals are to solve geometric problems in 2D, and to reach some possible generalizations in 3D.Some of the problems chosen will show the limits of DGS and CAS in the process of discovery and as well in the process of deductive proof. The problem of cutting a triangle in four equal parts will illustrate perfectly this dialectic. « Constructing from one given point of the plane two triangles of equal areas which bases are two given segments » is a problem that can enhance the use 3D DGS to investigate a possible generalization in3D. Eventually the Lhuillier problem will allow us to investigate in 3D with both CAS and DGS and state that these tools are only tools with their limits (the Lhuillier problem : given a first triangle, where are located points M of the plane which symmetric points with  respect to the sides of this triangle define a second triangle with the same area?).


\end{abstract_online}
    